\section{世界观：宿命论、自由意志与人生意义}

\subsection{确定性的马尔科夫过程}

在访谈的哲学部分，翁家翌展现了与技术话题截然不同的一面。他坚定地认为：

\begin{importantbox}{翁家翌的世界观：决定论}
\textbf{我们生活在一个确定性的马尔科夫过程里面。}所有的东西都是可以被预测的——你脑子里在想什么，下一个单词说什么，全都是宇宙大爆炸那一刻就定好了的。人没有自由意志，宏观世界不掷骰子（微观量子力学的随机性可以在"后台修改世界线"来解释）。
\end{importantbox}

他承认这是一个"相当悲观的世界观"，自己内心深处也不愿意接受——就好像自己变成了一个被模拟的原子。但基于他的个人经历，他认为这确实是事实。他表示"这个我已经验证无数遍了"，虽然不方便透露具体的个人经历。

主持人进一步追问：量子力学的不确定性呢？翁家翌的回答是"宏观不掷骰子，微观掷骰子"——但微观的随机性可以被解释为"在后台修改一些世界线"。他承认"你可以把我说的话认为是扯淡"，但他依然坚持这个观点。

这个世界观与他从事 RL 工作形成了有趣的张力：RL 的核心假设是 agent 可以通过选择不同 action 来改变未来的 reward，但如果世界是确定性的，那么 agent 的"选择"本身也是预先决定的。翁家翌是否意识到了这个张力？他的回答暗示了肯定——"有人想开发这种（预测未来的）模型，只是为了搞清楚这个世界背后运行的规律"。

\subsection{投资未来的悖论}

主持人敏锐地指出了翁家翌世界观中的矛盾：如果一切都是确定的，你投不投资未来都会到达那个点，那为什么还要投资？

翁家翌的回答是："投资一下还是会更好的嘛。"但当被追问"确定性的过程下你的投资不会有任何影响"时，他承认："投资未来可能也是确定性的——你投不投资也不是自己的自由意志。"

\subsection{人生迷茫期}

访谈结尾，翁家翌坦率地分享了自己当前的状态：

\begin{knowledgebox}{翁家翌的当下状态}
翁家翌处于职业生涯的某个迷茫期。曾经很喜欢的 RL Infra 工作，随着时间推移变得越来越确定性。他"曾经想通了自己想要什么，但现在又想不通了"。他的短期目标是提前退休，获得足够的资本，然后花时间找到自己真正想做的事。对未来10年的自己，他唯一的期望是"做自己那个时候想做的事，有足够的资源和能力"。
\end{knowledgebox}

\subsection{本章小结}

翁家翌的哲学思考为整个访谈增添了深度。一个搭建了 OpenAI 核心 infra 的工程师，在哲学上却是一个彻底的决定论者。他认为 AGI "板上钉钉"，剩下的都是"很确定性的事情"，"已经看到头了"。但这种确定性反而让他感到迷茫——如果一切都是确定的，努力的意义在哪里？

他对10年后的自己只有一个期望："做自己那个时候想做的事，有足够的资源和能力。"主持人追问"你不去干预他那时候想什么？"翁家翌回答："因为想法是会变的，可能你想什么也不重要。"然后他补充："我现在所能做的就是投资那个时候的我——还是投资未来，让他有选择的权利。"

他给出的最后留言是：\textbf{"去探索说到底自己想要什么——这个问题值得一生去思考。"}


